\section{Zusammenfassung und Fazit}
\label{section:summary}
\subsection{Zusammenfassung}

Das Ziel dieser Arbeit war es den vollständigen Entwurfsprozess einer Datenbank anhand eines selbstgewählten Beispiels - hier einer Carsharing-Applikation - zu erarbeiten. Im theoretischen Teil wurde der Entwurfsprozess an einem reduzierten Modell der Carsharing-Applikation von der anfänglichen Anforderungsanalyse, über den konzeptionellen Entwurf mittels Entity-Relationship-Modell, bis hin zum finalen DBMS (hier ein Relationenmodell) erläutert und gezeigt, welche Überlegungen und Methoden angewandt werden müssen, um die zentralen Ziele der Datenbankmodellierung -- geringe Redundanz, gute Bedienbarkeit, einfache Erweiterbarkeit, schneller Zugriff, Konsistenz, und Integrität -- schrittweise zu optimieren. Im praktischen Teil wurde der Modellierungsprozess anschließend auf die möglichst vollständig ausgeführte Carsharing-Applikation übertragen und mittels Racket und dem relationalen DBMS SQLite umgesetzt.

\subsection{Kritische Reflexion}

Bei der Ausarbeitung dieser Arbeit wurde deutlich, wie wichtig die anfänglich angewandten konzeptionellen Schritte im Datenbankentwurf sind, um ein relationales DBMS zu entwerfen und in einem realen DBMS umzusetzen. Angefangen bei der Anforderungsanalyse zeigte sich, dass die alleinige Perspektive aus Kundensicht -- wie sie beim Autor dieser Arbeit vorliegt -- nicht ausreicht um die Anforderungsanalyse vollständig abzubilden. Prozesse, die insbesondere die Mitarbeitenden der Buchhaltung oder auch den Service bei Reparaturen oder Kundenanfragen betreffen wurden teilweise gar nicht oder nur unvollständig abgebildet. Als Vorteil für die Ausarbeitung der Arbeit erwies sich die vergleichsweise geringe Komplexität des gewählten Beispiels im Vergleich zu einer realen Carsharing-Applikation. Diese geringe Komplexität erleichterte den konzeptionellen Entwurf im ERM als auch die Abbildung des ERMs in den Tabellen des relationalen DBMS. 

Bei der Datenbankmodellierung wurde trotz Durchführung der einzelnen Modellierungsschritte auch nach der Umsetzung des realen DBMS iterativ zwischen den einzelnen Modellierungsschritten gewechselt. Die verschiedenen Perspektiven der einzelnen Modellierungsebenen erleichterten die Übersicht über mögliche Schwachstellen oder Fehler im Datenbankdesign. Korrigiert wurden insbesondere die Entitätstypen und deren Beziehungstypen. Bei den Beziehungstypen lag die Schwierigkeit darin, sinnvolle, aber nicht unnötig definierte Fremdschlüsselattribute zu wählen. Es wurde deutlich, dass insbesondere für Einsteiger:innen eine gründliche Auseinandersetzung mit den Zielen der einzelnen Modellierungsschritte sehr hilfreich ist, um zeiteffizient und strukturiert ein reales DBMS umzusetzen. 

Schließlich ist noch zu erwähnen, dass wichtige Aspekte, die durch die Skalierung des DBMS auf eine reale Carsharing-Applikation, die Benutzerfreundlichkeit der Applikation, die intelligente Anbindung einer Fahrzeugflotte, als auch die Sicherheit hinsichtlich Daten- und Diebstahlschutz in dieser Ausarbeitung ausgelassen wurden.